% !TeX root = ../main.tex

\chapter{Improvements and Optimizations to Solution}
\label{chap:improvements}
Despite the fact that the completed system satisfied the acceptance criteria, some improvements are recommended for the future PCB revision.

The most important hardware improvement is the incorporation of the SDA polarity fix directly in the PCB design. The prototype required the SN74LVC2G04 inverter to correct the SDA signal polarity after the originally installed LTC2945 device variant turned out to be incompatible with the intended isolation interface. The future revision should use the correct LTC2945 variant throughout the design or incorporate the required inversion directly in the schematic and PCB.

All required I\textsuperscript{2}C pull-up resistors and defined logic states should also be incorporated directly in the next PCB revision. It will eliminate the prototype modification and improve the repeatability of assembly and startup.

The firmware initialization can also be made more tolerant of arbitrary power-up sequencing. During the development, the communication could be failed when the controller tried to initialize the sensor before the corresponding isolated rail circuitry was powered. The future firmware revision could periodically retry the unavailable sensors so that the communication is established automatically when the monitored rail is powered.

The per-channel measurement calibration is another potential improvement. The present system calculates the current using the configured current sense resistance. The measurement of the actual resistance or the determination of the calibration factor against the external reference will allow the individual gain correction to be stored for each channel. It will help to reduce the systematic errors caused by the shunt tolerance and allow tighter acceptance limits if the improved measurement accuracy is required.

The additional environmental characterization is also useful if the design is going to be developed beyond the capstone prototype. The testing at multiple ambient temperatures can be used to determine whether the temperature affects the measurement accuracy. EMC testing can determine the susceptibility to and generation of the electromagnetic interference. The long duration testing can be used to detect the intermittent communication, thermal, network or firmware problems that might not occur during the shorter acceptance test.

The automated test system should be retained for the future revisions. Since the raw measurements are saved separately from the report generation, the future hardware or firmware revisions can be tested using the same acceptance procedure and direct comparison between the revisions is possible.

These recommendations represent the refinements in reliability, manufacturability, calibration and qualification and not the changes to the basic architecture. The successful final test proves that the overall system concept is suitable for the further development.
